\section{Cranfield APA7 Author-Date Referencing}
\label{sec:authordaterefs}
Ensure the following lines are uncommented in the preamble to the template file
\begin{verbatim}
SHAUN ADD
\end{verbatim}
and towards the end of the file
\begin{verbatim}
SHAUN ADD
\end{verbatim}

We now cover a few of the examples from the APA7 Author-Date \cite{cranfield_library_service_apa7_2021}. 

\subsection{Journal Articles}
\label{sec:NumberedJournal}
\begin{description}
\item[Journal article with one author and no DOI]:
\begin{description}
\item[Parenthetical citation:] A recent article discussed the fracture... \parencite{hutchinson_fracture_2010}.
\item[Narrative citation:] Recently, \textcite{hutchinson_fracture_2010} discussed the fracture ...
\end{description}

\item[Journal article with two authors]:
\begin{description}
\item[Parenthetical citation:] It was not possible to separate targets \parencite{yang_perceptual_2016}.
\item[Narrative citation:]  \textcite{yang_perceptual_2016} found it was not possible to separate targets.
\end{description}

\item[Journal article with three to twenty authors]:
\begin{description}
\item[Parenthetical citation:] Previous authors \parencite{langley_process_2013} have identified assumptions \ldots .
\item[Narrative citation:]  \textcite{langley_process_2013} identify assumptions\ldots .
\end{description}

\item[Journal article with twenty one or more authors]:
If there are twenty-one or more authors, in the reference list add the names of the first nineteen
authors followed by \ldots and then the last author’s name.

\item[Online journal article with DOI] The synthesis of a new polymer of
cadmium(II) \linebreak ethylenediamineazide was determined by \ldots \parencite{yang_preparation_2010}. 
\item[Online journal article with URL] The synthesis of a new polymer of
cadmium(II) \linebreak ethylenediamineazide was determined by \ldots \parencite{yang_another_2010}.
\item[Journal article with no author]  TC-32B QT was demonstrated at an
exhibition held in Kent \cite{anon_first_2005}.
\end{description}

\subsection{Newspaper Articles}
\label{Sec:Newspapers}
If you are citing an online news website, e.g., BBC News, format the reference as a webpage (refer
to the section on the Internet).

\begin{description}
\item[Printed newspaper article] A blast at an atomic weapons factory was
met with denials by Iran \cite{frenkel_europe_nodate}.\\
\textbf{Note:} Use bibtex \texttt{misc} type, omit date from \texttt{year} field and place date in \texttt{note} field. 
\item[Article from an online newspaper]  Rayner \cite{rayner_combat_2011} reports that The Royal British Legion Centre for Blast Injury Studies is
to design…
\end{description}

\subsection{Books}
\label{sec:NumberedBooks}
\begin{description}
\item[Book with a single author] According to \textcite[p.35]{sommerville_software_2007} the most important part of software engineering is…
\item[Book with two authors] \textcite[p.65]{stair_principles_2010} suggest…
\item[Book with four or more authors] This was noted by \textcite[p.15]{payne_relationship_1995}
\item[Book with an editor] Latest thinking on gender in organizations
suggests… \cite{kumra_oxford_2014}
\item[Chapter of an edited book] \textcite[p.35]{haase_standards_2005} identified...
\item[Chapter of an edited book with organisation as author e.g. handbooks] 
The steel corrosion showed …\parencite{asm_international_corrosion_1998}
\item[Book with no author] In order to vote in the UK…\parencite{noauthor_whitakers_2011}
\item[Electronic book (eBook) or an online book] The project manager’s role can be defined as… \parencite{kerzer_project_2009}
\item[Book accessed via an eReader e.g. Kindle, iPad etc.]
Brodie tried to decipher the \linebreak code~\cite{dennis_secret_2012}\\
\textbf{NOTE: Place eReader type in square brackets at the end of the bibtex title field}
\item[Translated Book] In his discussion about algebraic
topology, \textcite[p.17]{matvev_lectures_2006} considered...\\
\textbf{Translator not yet supported}
\end{description}



\subsection{Conference Papers}
\label{sec:NumberedConference}
\textbf{NOTE:} Use bibtex entry \texttt{incollection}.
\begin{description}
\item[Full conference proceeding] The IEEE \cite[p.59]{noauthor_11_2010} suggests that \ldots
\item[Individual paper] \textcite{moon_aeroacoustic_2003} provide computations \ldots\\
\textbf{NOTE:} Place paper number in the bibtex \texttt{number} field.
\item[Online proceeding or paper] \textcite{wessel_assessing_2010} provide a ten step method\ldots
\item[Offshore Technology Conference (OTC) paper] According to \textcite{chiasson_large_1999} large bore \ldots
\end{description}

\subsection{Internet}
\label{sec:Internet}
\textbf{NOTES:} 
\begin{itemize}
\item The bibtex style file does not support the \texttt{accessed} attribute so include it in the note field of the bibtex entry.
\item If using Zotero then use a \texttt{document} and not a \texttt{web} entry as the \texttt{publisher} field is often used
\end{itemize}
\begin{description}
\item[Web page with an individual author] \textcite{bradley_notitle_2011} states that libraries and \ldots
\item[Web page with an organisation as author] The Natural Hazards Centre \cite{university_of_colorado_boulder_natural_2011} advised
that \ldots
\item[Blog] \textcite{borger_afghanistan_2010} claims in `Afghanistan
options are running out' \ldots\\
\textbf{NOTE: add [Blog] to end of title field.}
\item[Wiki] The CILIP ARLG National committee
meeting notes show \ldots \cite{barefoot_october_2012}
\textbf{NOTE: add [Wiki] to end of title field.}
\end{description}

\subsection{Thesis or dissertation}
\label{sec:thesis}
\begin{description}
\item[Print thesis or dissertation] An examination of Forth’s \citeyear{forth_morphological_1989}
research shows that…
\item[Online thesis or dissertation] \textcite{tatham_confidence_2010} identifies three areas
undergoing notable change…\\
\textbf{NOTE: use of question mark in thesis title is not presently detected in bibtex style file.}
\end{description}


\subsection{Course materials}
\label{sec:course}
\begin{description}
\item[Lecture or seminar] The lecture provoked much debate~\cite{johnson_peat_2012}
\item[Lecturers’ notes on Virtual Learning Environment \eg\ Canvas] As illustrated in the course materials~\cite{forth_analysis_2012}
\item[Course Virtual Learning Environment discussion forum] The debate during the lecture continued in the discussion forum. \textcite{hurst_are_2012} was
particularly vociferous.
\end{description}

\subsection{Official Publications}
\label{sec:officialPubs}
\textbf{Note:} These have not been tested yet but can probably be achieved with a bibtex \verb#misc# entry.
\begin{description}
\item[UK Statute (Act of Parliament)] 
\item[House of Commons or House of Lords Paper]
\item[House of Commons or House of Lords Bill]
\item[Statutory Instruments (SIs)]
\item[Hansard Debate]
\item[Command Paper (Green or White Paper)]
\item[European Directive]
\end{description}

\subsection{Reports}
\label{sec:Reports}
\begin{description}
\item[Report] An analysis of China’s investment in companies in different countries has shown... \cite{wolf_chinas_2011}.
\item[Online report] The effect was… \cite{curry_-flight_1989}.\\
\textbf{Note:} Presently unable to get colon (:) between report number and report title.
\item[Online industry or market research report] Media trends in the UK…\cite{marketline_media_2013}
\item[Online company or annual report] Despite economic uncertainty the
company’s profits did increase \cite{rolls-royce_group_plc_annual_2011}.
\item[Online company or financial report or data source] Rolls-Royce plc turnover has dramatically
reduced in 2009-2010 \cite{bureau_van_dijk_rolls-royce_2010}.
\end{description}


\subsection{Case Studies}
\textbf{Note:} These have not been tested yet but can probably be achieved with a bibtex \verb#misc# entry.

\subsection{Working Papers}
\textbf{Note:} These have not been tested yet but can probably be achieved with a bibtex \verb#misc# entry.




\subsection{Standards, patents, protocols and datasheets}
\label{sec:standards}
\textbf{Note:} 
\begin{itemize}
\item Use bibtex \texttt{report} item
\item Presently unable to get colon (:) between report number and report title.
\end{itemize}
\begin{description}
\item[Standard] In projects "specification, schedule and
cost always have to be bounded to ensure
ongoing viability" \cite{british_standards_institution_project_2010}.
\item[Online standard from a database] Some standards provide guidelines for effective project management \cite{british_standards_institution_project_2010}.
\item[Patent] The method of assembly... \cite{hodge_rear_2008}
\item[Joint Service Publication or protocol] The MOD has a simple definition of knowledge \cite{great_britain_ministry_of_defence_retaining_2011}.\\
\textbf{Note:} As the actual url is not available, put the availability and accessed data in the bibtex \texttt{note} field.
\item[Datasheet \eg ESDU] The buckling of the panel… \cite{esdu_elastic_2008}.
\end{description}

\subsection{Datasets and software}
\label{sec:datasets}
\begin{itemize}
\item Use bibtex \texttt{misc} type.
\item Include version number in parenthesis in the \texttt{title} field.
\end{itemize}
\begin{description}
\item[Dataset] Data from \textcite{partridge_spectra_2014} showed \ldots
\item[Software] Plone is an example of this sort of use
of an open source content management
system~\cite{plone_foundation_plone_2016}.
\end{description}

\subsection{Audio visual materials}
\textbf{Note:} These have not been tested yet but can probably be achieved with a bibtex \verb#misc# entry.

\subsection{Images}
\label{sec:images}
\textbf{Notes:} 
\begin{itemize}
\item Use \texttt{misc} type
\item For items available online   so include the \texttt{accessed} attribute in the \texttt{note} field of the bibtex entry.
\end{itemize}
\begin{description}
\item[Illustration, figure, diagram, logo or table] The estimated sales forecast is shown in \reffig{fig:sales}.
\begin{figure}[ht]
\begin{center}
\includegraphics[width=0.6\textwidth]{Sales}
\end{center}
\caption{Estimated Sales Forecast~\parencite[p.4]{wisker_good_2012}}
\label{fig:sales}
\end{figure}
\item[Embedded illustration, cartoon, diagram, etc.] The pressure on the housing market… \cite{wood_cartoon_2014}.\\
\textbf{Note:} Include the \texttt{[medium]} , here a \texttt{[cartoon]} at the start of the bibtex title field.
\item[Ordnance Survey map] P.O. equates to Post Office \parencite{ordnance_survey_oxford_1994}.\\
\textbf{Note:} 
\begin{itemize}
\item Use the \texttt{title} field to add the map number and scale.
\item Use the \texttt{note} field to add the map series.
\end{itemize}
\item[Online map] There are many farms in the area \parencite{ordnance_survey_peak_2010}.\\
\textbf{Note:} The bibtex style file does not support the \texttt{accessed} attribute so include it in the note field of the bibtex entry.
\item[Painting or drawing] \textbf{Note:} This has not been tested yet but can probably be achieved with a bibtex \verb#misc# entry.
\item[Painting or drawing online] Groves’ {\em Carrot on a stick} \citeyear{groves_carrot_2007} is playing with colours. 
\item[Photograph – print or slide] 
Sublime lighting captures the atmosphere
wonderfully \cite{collins_pictures_2010}\\
\textbf{Note:} I uses a \texttt{book} entry here.
\item[Photograph online or in an online collection \eg Flickr, Instagram] An abundance of movement is
captured \cite{rush_running_1997}.
\end{description}


\subsection{Personal communications}
\textbf{Notes:}\\
\begin{itemize}
\item Use \texttt{misc} type
\item If dates are not fully listed then place them in the \texttt{note} field.
\end{itemize}
\begin{description}
\item[Email or letter] The details of the organisational structure
were outlined \cite{godfrey_restructure_nodate}.
\item[Conversation] It was felt the organisational culture
changed when the new CEO arrived \cite{williams_conversation_nodate}.
\item[Interview, survey or questionnaire] The theory was dismissed as
nothing new \cite{benson_should_nodate}.
\end{description}

